\documentclass[11pt]{article}
\usepackage[T1]{fontenc}
\usepackage[margin=1in]{geometry}
\usepackage{amsmath,amssymb}
\usepackage[hidelinks]{hyperref}

\title{Literature Status: Dimension-Free Convexity of Fermionic Entropy}
\author{Automatic Math Discovery Run \texttt{fa\_banach\_001}, lane 0}
\date{August 9, 2026}

\begin{document}
\maketitle

\section*{Status}

This is a \textbf{literature-already-answered} packet. It records a full
positive resolution, already present in later literature, of the
all-dimensional Hessian and geodesic-convexity conjecture from
arXiv:1203.5377. It is not a new proof packet.

\section*{Source Question}

Carlen and Maas construct a non-commutative Riemannian transport metric for
the fermionic Fokker--Planck equation in arXiv:1203.5377 \cite{CM14}.
Remark~5.6, on PDF page 24, proves for Clifford dimensions \(n=1,2\) that
\[
  \operatorname{Hess}_{\rho} S(\nabla U,\nabla U)
  \geq \|\nabla U\|_{\rho}^{2},
\]
and conjectures that this inequality holds for every \(n\).

\section*{Supporting Answer}

The introduction of arXiv:1609.01254, on PDF page 2, explicitly says that the
paper proves the sharp geodesic-convexity result for the Fermi
Ornstein--Uhlenbeck semigroup and thereby solves the problem left open in
\cite{CM14}. This satisfies the explicit-awareness criterion for the
literature-already-answered classification.

The decisive statement is Theorem~8.6 on PDF page 34 \cite{CM17}. For inverse
temperature \(\beta\geq0\), it proves that relative entropy is geodesically
\(\lambda_\beta\)-convex in the Riemannian metric associated with the Fermi
Ornstein--Uhlenbeck generator, where
\[
  \lambda_\beta=\min_j\cosh(\beta e_j/2).
\]
The paragraph immediately preceding that theorem explains that the metric can
be written using the skew derivations exactly as in \cite{CM14}. At
\(\beta=0\), the invariant state is tracial and \(\lambda_0=1\). Thus the
theorem gives geodesic \(1\)-convexity for the source metric in every Clifford
dimension. In this smooth finite-dimensional Riemannian setting, this is
equivalent to
\[
  \operatorname{Hess}_{\rho} S \geq g_{\rho},
\]
which is precisely the conjectured inequality.

\section*{Scope}

The Hessian and geodesic-convexity conjecture is fully answered positively,
with the sharp constant \(1\). The source paper's separate Remark~5.7 question
about the sectional curvature of the density manifold is not answered by this
identification.

\section*{Search Evidence}

The run's cheap indexes were searched for arXiv:1203.5377 and the terms
``fermionic'', ``Hessian'', ``entropy convexity'', ``Ricci'', and ``Clifford'';
no prior exact packet was present. A bounded local-source and arXiv search
located arXiv:1609.01254. Its introduction explicitly identifies the earlier
problem as solved, and Theorem~8.6 supplies the exact all-dimensional sharp
statement.

\begin{thebibliography}{9}

\bibitem{CM14}
E.~A. Carlen and J. Maas,
\emph{An Analog of the 2-Wasserstein Metric in Non-commutative Probability
under which the Fermionic Fokker-Planck Equation is Gradient Flow for the
Entropy}, arXiv:1203.5377; Comm. Math. Phys. 331 (2014), 887--926.

\bibitem{CM17}
E.~A. Carlen and J. Maas,
\emph{Gradient Flow and Entropy Inequalities for Quantum Markov Semigroups
with Detailed Balance}, arXiv:1609.01254; J. Funct. Anal. 273 (2017),
1810--1869.

\end{thebibliography}

\end{document}
